\chapter{Empirical Model Benchmarking and Latency Profiling}
\label{ch:appendix_evaluation_benchmarks}

\section{Overview of Empirical Benchmark Charts}
\label{sec:appendix_benchmarks_overview}
The following figures in this appendix are the graph and chart representations of the experimental results analysis covered in Chapter~\ref{ch:results}. They include comparisons among deep learning methods, confusion matrices for each finger, accuracy of AprilTag homography tracking in terms of tilt angles, and runtime CPU latency.

\section{Model Architecture Accuracy versus CPU Latency}
\label{sec:appendix_model_benchmark_chart}
The Macro F1-score and CPU inference time of various deep learning architectures have been compared in Figure~\ref{fig:appendix_benchmark_acc} of the current benchmark test. The uni-directional LSTM with scale-normalized position and velocity scores best among all architectures with an F1-score of 0.963 and CPU inference time of $2.08\text{ ms}$ on the CPU.

\begin{figure}[htbp]
\centering
\includegraphics[width=0.88\textwidth]{figures/fig_model_benchmark_acc.png}
\caption{Comparison of Macro F1-score and CPU inference latency across evaluated deep learning model architecture families.}
\label{fig:appendix_benchmark_acc}
\end{figure}

\newpage

\section{Per-Digit Touch Confusion Matrix}
\label{sec:appendix_confusion_matrix_chart}
As shown in Figure~\ref{fig:appendix_confusion_matrix}, the normalized confusion matrix for our chosen PyTorch LSTM model was generated using 5,000 test windows. All five fingers have touch detection accuracies greater than 93\%. The highest accuracy belongs to the index finger (99.0\% recall), while the pinky finger is slightly lower due to its smaller movement.

\begin{figure}[htbp]
\centering
\includegraphics[width=0.82\textwidth]{figures/fig_confusion_matrix.png}
\caption{Normalized confusion matrix of the PyTorch LSTM touch classifier across all five individual finger digits.}
\label{fig:appendix_confusion_matrix}
\end{figure}

\newpage

\section{AprilTag Homography Tracking Under Perspective Tilt}
\label{sec:appendix_homography_tilt_chart}
The influence of camera tilt angles (ranging from $0^\circ$ to $75^\circ$) on AprilTag tracking is presented in Figure~\ref{fig:appendix_homography_tilt}. The corner detection percentage remains close to 100\%, while the average tracking error remains at less than $0.42\text{ mm}$ for angles ranging from $0^\circ$ to $60^\circ$ of the desk. Beyond $65^\circ$, there is an increase in error due to perspective tilt, showing that $0^\circ$ to $60^\circ$ is the best angle range.

\begin{figure}[htbp]
\centering
\includegraphics[width=0.88\textwidth]{figures/fig_homography_tilt_error.png}
\caption{AprilTag corner detection rate and RMS planar reprojection error across camera perspective tilt angles.}
\label{fig:appendix_homography_tilt}
\end{figure}

\newpage

\section{Multi-Threaded Real-Time Pipeline Latency Breakdown}
\label{sec:appendix_latency_breakdown_chart}
As can be seen in Figure~\ref{fig:appendix_latency_breakdown}, this represents the time required for each pipeline stage to execute itself on an Intel Core i7 processor. Total time required for processing is $29.09\text{ ms}$ (approximately $34.4\text{ FPS}$). It utilizes only about 35\% of the total possible $83.33\text{ ms}$ budget at 12~FPS, which means the program runs smoothly in real time without needing a GPU.

\begin{figure}[htbp]
\centering
\includegraphics[width=0.88\textwidth]{figures/fig_latency_breakdown.png}
\caption{End-to-end CPU execution latency breakdown per frame across the multi-threaded virtual keyboard pipeline.}
\label{fig:appendix_latency_breakdown}
\end{figure}
